\section{The coordinate integrals in the Jacquet estimates}
\label{jk-section}

This section proves the coordinate estimates used in the
construction of the integrable majorant for the Jacquet pairing.
All the integrals are Lebesgue integrals of nonnegative functions.
They are therefore well defined, possibly with value \(+\infty\),
before convergence is established.  Every interchange below follows
from Tonelli's theorem.  We give an exact value and a necessary and
sufficient convergence criterion for the basic integral; the
weighted estimate then follows from elementary averaging.
The notation and the precise correspondence with the source are
recorded at the end of the section.

\subsection{The basic integral and its exact convergence region}

Let \(p,m\) be nonnegative integers.  On
\[
 X_{p,m}=\mathbb C^p\times\mathbb C^p\times
                  \mathbb C^m\times\mathbb R
\]
write \(x=(a,b,z,t)\), with indices \(0\leq j<p\).
Lebesgue measure on \(\mathbb C\) means
\(d(\operatorname{Re}z)\,d(\operatorname{Im}z)\);
in dimension zero it has total mass one.
For \(0\leq l\leq p\), put
\begin{align}
 L_l(x)&=1+\sum_{j<l}|a_j|^2,\nonumber\\
 B_l(x)&=-\tfrac12\|z\|^2+it-\sum_{j\geq l}b_j\overline{a_j},
       \label{jk-coordinates}\\
 R_l(x)&=L_p(x)+\|z\|^2+\sum_{j\geq l}|b_j|^2,\nonumber\\
 d_l(x)&=\bigl(L_l(x)R_l(x)+|B_l(x)|^2\bigr)^{1/2}.
       \nonumber
\end{align}
In particular \(d_l\geq1\).  Define
\[
 \sigma_l=\begin{cases}2&l<p,\\m+1&l=p,\end{cases}
 \qquad
 I_{p,m}(\nu)=
       \int_{X_{p,m}}\prod_{l=0}^{p}d_l(x)^{-\sigma_l-\nu_l}\,dx,
       \qquad \nu\in\mathbb R^{p+1}.
\]
Write \(S_l=\nu_0+\cdots+\nu_l\).

\begin{theorem}\label{jk-exact}
The integral \(I_{p,m}(\nu)\) is finite if and only if
\begin{equation}\label{jk-cone}
 S_l>0\quad(0\leq l\leq p),\qquad
 S_l+2\sum_{j=l+1}^{p}\nu_j>0\quad(0\leq l<p).
\end{equation}
On this region it equals
\begin{equation}\label{jk-product}
 I_{p,m}(\nu)=
 \mathcal B_m(S_p)
 \prod_{l=0}^{p-1}
   \frac{(2\pi)^2}{S_l\bigl(S_l+2\sum_{j>l}\nu_j\bigr)},
\end{equation}
where, for \(v>0\),
\begin{equation}\label{jk-base-value}
 \mathcal B_m(v)=
  2^m\pi^{m+1/2}
  \frac{\Gamma((m+v)/2)\Gamma(v)}
       {\Gamma((m+1+v)/2)\Gamma(m+v)} .
\end{equation}
An empty product has value one.
\end{theorem}

\begin{proof}
First take \(p=0\).  Directly from \eqref{jk-coordinates},
\[
 d_0^2=(1+\|z\|^2/2)^2+t^2.
\]
For \(q>1/2\) and \(A>0\), integration of
\[
 u^{-q}=\frac{1}{\Gamma(q)}
                   \int_0^\infty s^{q-1}e^{-su}\,ds
\]
against \(u=A^2+t^2\) gives
\[
 \int_{\mathbb R}(A^2+t^2)^{-q}\,dt
  =\sqrt\pi\,\frac{\Gamma(q-\tfrac12)}{\Gamma(q)}
                      A^{1-2q}.
\]
Here we used \(\int_{\mathbb R}e^{-st^2}\,dt=\sqrt{\pi/s}\).
The same calculation in \(m\) complex coordinates gives, for \(h>m\),
\[
 \int_{\mathbb C^m}(1+\|z\|^2/2)^{-h}\,dz
                    =(2\pi)^m\frac{\Gamma(h-m)}{\Gamma(h)}.
\]
These identities follow directly from the displayed Laplace
integral and Tonelli; the Gaussian identity follows by squaring the
one dimensional integral and using polar coordinates in
\(\mathbb R^2\).
Apply them with \(q=(m+1+v)/2\) and \(h=m+v\).
This proves \eqref{jk-base-value} when \(v>0\).
If \(m+v\leq0\), the integral in \(t\) already diverges.
If \(m+v>0\) but \(v\leq0\), integration in \(t\) leaves
\((1+\|z\|^2/2)^{-m-v}\), whose integral diverges:
for \(m\geq1\), polar coordinates give a tail comparable to
\(\int_1^\infty r^{-1-2v}\,dr\).
For \(m=0\) the first alternative applies.
Thus the base integral is finite exactly when \(v>0\).

Now suppose \(p\geq1\), and set \(L=1+|a_0|^2\).
Completing the square gives the identity
\begin{equation}\label{jk-square}
 d_0^2=
 L\left|b_0-\frac{a_0B_1}{L}\right|^2+\frac{d_1^2}{L}.
\end{equation}
Indeed \(B_0=B_1-b_0\overline{a_0}\),
\(R_0=R_1+|b_0|^2\), and \(L_0=1\);
expansion of both sides gives the same expression.
The remaining \(d_l\), \(l\geq1\), do not depend on \(b_0\).
For fixed remaining coordinates, integration in \(b_0\) is finite
exactly when \(\nu_0>0\).  In that case polar coordinates give
\begin{equation}\label{jk-first-integration}
 \int_{\mathbb C}d_0^{-2-\nu_0}\,db_0
       =\frac{2\pi}{\nu_0}L^{\nu_0/2-1}d_1^{-\nu_0}.
\end{equation}
For \(\nu_0\leq0\), the two dimensional tail diverges, independently
of all the fixed coordinates.  The full integral is then infinite.

For fixed \(a_0\), make the substitution
\[
 a_j=\sqrt L\,\widetilde a_{j-1},\quad
 b_j=\sqrt L\,\widetilde b_{j-1}\quad(j\geq1),\qquad
 z=\sqrt L\,\widetilde z,\quad t=L\widetilde t.
\]
Its Jacobian is \(L^{2(p-1)+m+1}\).  The definitions show that
\[
 L_{l+1}=L\widetilde L_l,\quad
 B_{l+1}=L\widetilde B_l,\quad
 R_{l+1}=L\widetilde R_l,\quad
 d_{l+1}=L\widetilde d_l.
\]
The sum of the remaining critical exponents is
\(2(p-1)+m+1\), exactly the Jacobian exponent.
Consequently the remaining factor in \(a_0\) is
\[
 (1+|a_0|^2)^{-1-\nu_0/2-\sum_{j\geq1}\nu_j}.
\]
Its integral is finite exactly when
\(\nu_0+2\sum_{j\geq1}\nu_j>0\), and then equals
\(2\pi/(\nu_0+2\sum_{j\geq1}\nu_j)\).
Thus, whenever these two inequalities hold,
\begin{equation}\label{jk-recursion}
 I_{p,m}(\nu)=
 \frac{(2\pi)^2}
      {\nu_0(\nu_0+2\sum_{j\geq1}\nu_j)}
 I_{p-1,m}(\nu_0+\nu_1,\nu_2,\ldots,\nu_p).
\end{equation}
This is also an identity of extended nonnegative integrals.
The lower dimensional integral is strictly positive, so divergence
of either factor cannot be concealed by a zero factor.
Induction now proves necessity and sufficiency of all the
inequalities in \eqref{jk-cone}, and multiplication of the factors
in \eqref{jk-recursion} proves \eqref{jk-product}.
\end{proof}

\begin{corollary}\label{jk-uniform}
On every compact subset of the open region \eqref{jk-cone},
the integrands defining \(I_{p,m}\) have a common integrable
majorant.  In particular the integral depends continuously on
\(\nu\) there.  The same assertion holds after any fixed number of
derivatives with respect to the exponents.
\end{corollary}

\begin{proof}
At any point \(\nu^*\) of the region, choose \(\varepsilon>0\)
such that \(\nu^*-2\varepsilon(1,\ldots,1)\) remains in the region.
For \(|\nu_l-\nu_l^*|<\varepsilon\), the inequality \(d_l\geq1\)
gives an integrable majorant with exponents
\(\nu_l^*-\varepsilon\).
For derivatives, each additional factor is a product of powers
of \(\log d_l\).  The elementary bound
\((\log u)^N\leq C_{N,\varepsilon}u^\varepsilon\), \(u\geq1\),
gives the majorant with exponents
\(\nu_l^*-2\varepsilon\), with a smaller \(\varepsilon\) if needed.
A finite covering of the compact set supplies a finite sum of
such majorants.  Dominated convergence proves the assertions.
\end{proof}

\subsection{An integral with one coordinate removed}

For \(a\in\mathbb C^{p+1}\), \(b\in\mathbb C^p\),
\(z\in\mathbb C^m\), \(t\in\mathbb R\), regard
\[
 x=(a,(0,b_0,\ldots,b_{p-1}),z,t)\in X_{p+1,m}.
\]
Use the quantities \(L_l,d_l,\sigma_l\) for this larger space.
For \(\alpha\in\mathbb R\) and \(\nu\in\mathbb R^{p+1}\), define
\[
 T_{p,m}(\alpha,\nu)=
 \int L_1^{(\alpha-2)/2}d_1^{-\alpha}
                \prod_{i=0}^{p}d_{i+1}^{-\sigma_{i+1}-\nu_i}
                                      \,da\,db\,dz\,dt .
\]

\begin{corollary}\label{jk-tail}
This integral is finite exactly when
\[
 \alpha+2\sum_{i=0}^{p}\nu_i>0
 \quad\hbox{and}\quad
 (\nu_0+\alpha,\nu_1,\ldots,\nu_p)
                 \text{ satisfies \eqref{jk-cone}}.
\]
Under these conditions,
\begin{equation}\label{jk-tail-value}
 T_{p,m}(\alpha,\nu)=
       \frac{2\pi}{\alpha+2\sum_i\nu_i}
          I_{p,m}(\nu_0+\alpha,\nu_1,\ldots,\nu_p).
\end{equation}
No sign condition on \(\alpha\) is required.
\end{corollary}

\begin{proof}
Set \(L=1+|a_0|^2\) and use the same scaling as in the
proof of Theorem~\ref{jk-exact}.
The Jacobian is now \(L^{2p+m+1}\) and cancels the sum
of the critical exponents of \(d_1,\ldots,d_{p+1}\).
The residual power is
\(L^{-1-\alpha/2-\sum_i\nu_i}\).
Its integral over \(a_0\in\mathbb C\) gives the displayed factor
and its exact convergence condition.  The other factor is the
strictly positive integral on the right of \eqref{jk-tail-value}.
Tonelli proves both directions, including divergence cases.
\end{proof}

\subsection{Three elementary estimates}

\begin{lemma}\label{jk-sublevel}
Let \(P(u)=c_0+c_1u+c_2u^2\), with complex coefficients, and let
\(M=\max_j|c_j|\).  If \(M>0\) and \(\varepsilon>0\), then
\begin{equation}\label{jk-sublevel-measure}
 \bigl|\{u\in[1,2]:|P(u)|\leq\varepsilon M\}\bigr|
                                  \leq \sqrt{432\varepsilon}.
\end{equation}
Consequently, for \(0<\delta\leq1/2\),
\begin{equation}\label{jk-sublevel-integral}
 \int_1^2(1+|P(u)|)^{-\delta}\,du
                         \leq30(1+M)^{-\delta/2}.
\end{equation}
\end{lemma}

\begin{proof}
Let \(E\) be the set in \eqref{jk-sublevel-measure} and
\(\mu=|E|>0\).  Select three points \(u_1,u_2,u_3\) of \(E\)
at pairwise distance at least \(\mu/6\): after selecting each
of the first two points, delete its interval of radius
\(\mu/6\); the deleted measure is at most \(2\mu/3\).
Lagrange interpolation gives
\[
 P(u)=\sum_{i=1}^{3}P(u_i)
              \prod_{j\ne i}\frac{u-u_j}{u_i-u_j}.
\]
Each coefficient of the numerator \(\prod_{j\ne i}(u-u_j)\)
has absolute value at most \(4\), since \(1\leq u_j\leq2\).
Each denominator has absolute value at least \(\mu^2/36\).
Comparing coefficients yields
\(M\leq432\varepsilon M/\mu^2\), proving
\eqref{jk-sublevel-measure}.

For \(M\geq1\), split at \(|P|=\sqrt M\).
The contribution from the first set is at most
\(\sqrt{432}M^{-1/4}\); on its complement the integrand
is at most \(M^{-\delta/2}\).
Since \(\delta\leq1/2\), their sum is at most
\((\sqrt{432}+1)M^{-\delta/2}
 \leq30(1+M)^{-\delta/2}\).
For \(M\leq1\), including \(M=0\), the integral is at most one,
which satisfies the same bound.
\end{proof}

\begin{lemma}\label{jk-riesz}
If \(0<e<\gamma<2\), there is a finite constant
\(C_{e,\gamma}\), independent of \(c\in\mathbb C\) and \(D>0\),
such that
\begin{equation}\label{jk-riesz-bound}
 \int_{\mathbb C}
 (|w-c|^2+D^2)^{-1+e/2}|w|^{-\gamma}\,dw
       \leq C_{e,\gamma}(|c|^2+D^2)^{-(\gamma-e)/2}.
\end{equation}
\end{lemma}

\begin{proof}
Scale by \(r=(|c|^2+D^2)^{1/2}\).  The resulting factor is
\(r^{e-\gamma}\), and the new parameters satisfy
\(|c|^2+D^2=1\).
For \(|w|\geq3\), the integrand is at most a constant times
\(|w|^{-2+e-\gamma}\); its integral is uniformly finite because
\(\gamma>e\).
It remains to consider \(|w|\leq3\).
If \(|c|\leq1/2\), then \(D\geq\sqrt3/2\), and the only
possible singularity is bounded by a constant times
\(|w|^{-\gamma}\), integrable because \(\gamma<2\).
If \(|c|>1/2\), the discs \(|w|<1/4\) and
\(|w-c|<1/4\) are disjoint.
On the first, the other factor is bounded uniformly.
On the second, \(|w|^{-\gamma}\) is bounded uniformly, while
\[
 (|w-c|^2+D^2)^{-1+e/2}\leq |w-c|^{-2+e},
 \qquad w\ne c,
\]
which is integrable because \(e>0\).
On the remainder both distances are at least \(1/4\), so
the integrand is bounded uniformly.
This proves the normalized estimate and hence the result.
\end{proof}

\begin{lemma}\label{jk-slow}
Let \(0<\kappa<1\), \(R>0\), and
\[
 w(\beta)=\prod_{i=1}^{n}
                (A_i^2+(\beta-b_i)^2)^{-t_i/2},
 \qquad
 A_i\geq R,\quad t_i\geq0,\quad \sum_i t_i\leq T .
\]
For every \(c\in\mathbb R\),
\begin{equation}\label{jk-slow-bound}
 \int_{\mathbb R}w(\beta)|\beta-c|^{-\kappa}\,d\beta
 \leq
 \left(1+\frac{2\,19^{\max(T,0)/2}}{1-\kappa}\right)
 R^{-\kappa}\int_{\mathbb R}w(\beta)\,d\beta .
\end{equation}
The assertion is valid for extended nonnegative integrals.
\end{lemma}

\begin{proof}
Translate to \(c=0\).  The case \(\int w=+\infty\) is
immediate, so suppose that integral is finite.
For \(|\beta|\leq R\) and \(R\leq s\leq2R\),
\[
 A_i^2+(s-b_i)^2
 \leq A_i^2+2(\beta-b_i)^2+18R^2
 \leq19\bigl(A_i^2+(\beta-b_i)^2\bigr).
\]
Therefore \(w(\beta)\leq19^{\max(T,0)/2}w(s)\).
Averaging over \(s\in[R,2R]\) gives
\[
 w(\beta)\leq19^{\max(T,0)/2}R^{-1}
                       \int_{\mathbb R}w(s)\,ds
             \qquad (|\beta|\leq R).
\]
Integration of \(|\beta|^{-\kappa}\) on \([-R,R]\) gives
the second term in the constant.
On \(|\beta|\geq R\), use
\(|\beta|^{-\kappa}\leq R^{-\kappa}\).
\end{proof}

\subsection{A finite integral below the critical exponents}

\begin{proposition}\label{jk-weighted}
Let \(P\geq1\), \(m\geq0\), \(0<\gamma\leq1\), and
\(e_0,\ldots,e_P>0\).  Suppose that
\begin{equation}\label{jk-loss-condition}
 e_0+2\sum_{l=1}^{P}e_l<\gamma/2.
\end{equation}
Use \eqref{jk-coordinates} on \(X_{P,m}\), and put
\(s_l=\sigma_l-e_l\).  Then
\begin{equation}\label{jk-weighted-integral}
 \int_{X_{P,m}}\prod_{l=0}^{P}d_l(x)^{-s_l}
                                  |b_0|^{-\gamma}\,dx<\infty .
\end{equation}
\end{proposition}

\begin{proof}
The hypotheses give \(0<e_0<\gamma<2\) and \(0<e_l<1\)
for all \(l\); hence \(s_l\geq0\).
Fix all variables except \(b_0\), and let \(L=1+|a_0|^2\).
Write \(D_0=d_0|_{b_0=0}\) and \(D_l=d_l\) for \(l\geq1\);
the latter quantities are independent of \(b_0\).
With \(c=a_0B_1/L\), identity \eqref{jk-square} gives
\[
 d_0^2=L|b_0-c|^2+D_1^2/L,\qquad
 D_0^2=L|c|^2+D_1^2/L .
\]
Applying Lemma~\ref{jk-riesz}, with its positive parameter
\(D=D_1/L\), yields
\begin{equation}\label{jk-b-integration}
 \int_{\mathbb C}d_0^{-2+e_0}|b_0|^{-\gamma}\,db_0
              \leq C L^{-1+\gamma/2}D_0^{-(\gamma-e_0)}.
\end{equation}
The powers of \(L\) here follow by factoring \(L\) from
\(d_0^2\) and using
\(|c|^2+(D_1/L)^2=D_0^2/L\).

Put \(\alpha=\gamma/2-e_0>0\) and \(\kappa=\gamma/2\).
For the point with \(b_0=0\), let
\[
 \theta_l=\operatorname{Im}\sum_{j\geq l}b_j\overline{a_j}.
\]
Then
\[
 D_l^2=A_l^2+(t-\theta_l)^2,\qquad
 A_l^2=L_lR_l+
       \left(-\|z\|^2/2-
                   \operatorname{Re}\sum_{j\geq l}
                                     b_j\overline{a_j}\right)^2.
\]
For \(l\geq1\), both \(L_l\) and \(R_l\) are at least \(L\);
thus \(A_l\geq L\).
Moreover \(B_0=B_1\) and \(R_0=R_1\) at \(b_0=0\).
It follows that
\[
 D_1\leq\sqrt L\,D_0,\qquad |t-\theta_1|\leq D_0.
\]
Apart from the single value \(t=\theta_1\), this implies
\[
 D_0^{-(\gamma-e_0)}
   \leq L^{\alpha/2}D_1^{-\alpha}|t-\theta_1|^{-\kappa}.
\]
Apply Lemma~\ref{jk-slow} to
\[
 w(t)=D_1^{-\alpha}\prod_{l=1}^{P}D_l^{-s_l}.
\]
Its exponents are nonnegative and have sum at most
\(\alpha+\sum_{l=1}^{P}\sigma_l\); all its parameters
\(A_l\) are at least \(L\).
The constant is therefore independent of the fixed coordinates.
After integrating in \(t\), the product of the powers of \(L\) is
\[
 L^{-1+\gamma/2}L^{\alpha/2}L^{-\kappa}
                                      =L^{(\alpha-2)/2}.
\]
By \eqref{jk-b-integration}, the integral in
\eqref{jk-weighted-integral} is consequently at most a constant
times
\[
 T_{P-1,m}\bigl(\alpha,(-e_1,\ldots,-e_P)\bigr).
\]
Condition \eqref{jk-loss-condition} says
\(\alpha>2\sum_{l=1}^{P}e_l\).
The first condition of Corollary~\ref{jk-tail} follows.
Every partial sum for the shifted vector
\((\alpha-e_1,-e_2,\ldots,-e_P)\) is positive.
Every other required expression is of the form
\[
 \alpha-\sum_{j\leq l}e_j-2\sum_{j>l}e_j
                       \geq\alpha-2\sum_{j=1}^{P}e_j>0,
\]
with the appropriate initial segment starting at \(j=1\).
Thus all conditions of that corollary hold, proving finiteness.
\end{proof}

\subsection{The estimate with a linear character}

For \(x=(a,b,z,t)\in X_{P,m}\), define the complex linear form
\[
 \ell(x)=\sum_{j=0}^{P-1}(a_j\eta_j+b_j\eta'_j)
                                      +\sum_{i=1}^{m}z_i\xi_i,
\]
where \(\eta,\eta'\in\mathbb C^P\), \(\xi\in\mathbb C^m\).

\begin{theorem}\label{jk-oscillatory}
Let \(P\geq1\), \(m\geq0\),
\[
 0<\delta\leq\tfrac12,\qquad
 0<\gamma\leq\delta/8,\qquad e_l>0,
 \qquad e_0+2\sum_{l=1}^{P}e_l<\gamma/2 .
\]
Set \(s_l=\sigma_l-e_l\).
There is a finite constant depending only on
\(P,m,\delta,\gamma,e\) such that, for every
\(\eta,\eta',\xi,c_0\) with \(\eta'_0\ne0\) and \(c_0\in\mathbb R\),
\begin{equation}\label{jk-os}
 \int_{X_{P,m}}\prod_{l=0}^{P}d_l(x)^{-s_l}
          \bigl(1+|2\operatorname{Im}\ell(x)+c_0|\bigr)^{-\delta}
             \,dx
 \leq C|\eta'_0|^{-\gamma}(1+|c_0|)^{-\delta/4}.
\end{equation}
In particular the bound is uniform in all the other coefficients
of the linear form.
\end{theorem}

\begin{proof}
Write \(F(x)=\prod_l d_l(x)^{-s_l}\).
We use two changes of variables preserving Lebesgue measure.
All the averaging is over \([1,2]\), an interval of length one.

First put
\[
 \rho(u)=\frac{(1+iu)^2}{1+u^2},\qquad
 \mathcal R_u(a,b,z,t)=(\rho(u)a,\rho(u)b,\rho(u)z,t).
\]
Since \(|\rho(u)|=1\), this map preserves measure.
All \(L_l,R_l,B_l\), and hence \(F\), are unchanged, whereas
\(\ell(\mathcal R_u x)=\rho(u)\ell(x)\).
For \(X\in\mathbb C\), direct multiplication gives
\[
 2\operatorname{Im}(\rho(u)X)+c_0
     =\frac{(2\operatorname{Im}X+c_0)
                 +4\operatorname{Re}X\,u
                 +(c_0-2\operatorname{Im}X)u^2}{1+u^2}.
\]
Let \(M\) be the maximum of the absolute values of the three
coefficients of the numerator.  Then
\[
 |c_0|\leq M,\qquad
 |\operatorname{Re}X|\leq M/4,\qquad
 |\operatorname{Im}X|\leq M/2,\qquad |X|\leq M.
\]
As \(1+u^2\leq5\), Lemma~\ref{jk-sublevel} shows that
\begin{equation}\label{jk-rotation-bound}
 \int_1^2
 (1+|2\operatorname{Im}(\rho(u)X)+c_0|)^{-\delta}\,du
 \leq30\,5^\delta
              (1+|c_0|)^{-\delta/4}(1+|X|)^{-\delta/4}.
\end{equation}
Indeed \((1+M)^2\geq(1+|c_0|)(1+|X|)\).
Changing variables by \(\mathcal R_u\), and then averaging,
bounds the left side of \eqref{jk-os} by
\begin{equation}\label{jk-after-rotation}
 30\,5^\delta(1+|c_0|)^{-\delta/4}
                 \int F(x)(1+|\ell(x)|)^{-\delta/4}\,dx .
\end{equation}

The second change of variables is
\[
 \mathcal S_u(a,b,z,t)=
 ((ua_0,a_1,\ldots),(b_0/u,b_1,\ldots),z,t).
\]
The real Jacobian is \(u^2u^{-2}=1\).
For \(1\leq u\leq2\), each \(L_l\) can only increase,
each \(R_l\) is at least one quarter of its original value,
and each \(B_l\) is unchanged.  Thus
\[
 d_l(\mathcal S_u x)\geq d_l(x)/2,\qquad
 F(\mathcal S_u x)\leq2^{\sum_l s_l}F(x).
\]
Here \(s_l\geq0\), as already observed in
Proposition~\ref{jk-weighted}.
Consequently change of variables and averaging give
\begin{align}
 \int F(x)(1+|\ell(x)|)^{-\delta/4}\,dx
 &\leq2^{\sum_l s_l}
   \int F(x)\int_1^2
       (1+|\ell(\mathcal S_u x)|)^{-\delta/4}\,du\,dx.
 \label{jk-after-scaling}
\end{align}
Write
\[
 A=a_0\eta_0,\qquad C=b_0\eta'_0,\qquad
 B=\sum_{j\geq1}(a_j\eta_j+b_j\eta'_j)+\sum_i z_i\xi_i.
\]
Then \(u\ell(\mathcal S_u x)=C+Bu+Au^2\).
Since \(u\leq2\), another application of
Lemma~\ref{jk-sublevel}, now with exponent \(\delta/4\), gives
\[
 \int_1^2(1+|\ell(\mathcal S_u x)|)^{-\delta/4}\,du
 \leq30\,2^{\delta/4}(1+|C|)^{-\delta/8}
 \leq30\,2^{\delta/4}|b_0\eta'_0|^{-\gamma}
\]
whenever \(b_0\eta'_0\ne0\).
The last inequality uses \(0<\gamma\leq\delta/8\).
As \(\eta'_0\ne0\), the excluded set \(b_0=0\) has measure
zero.  Alternatively the bound holds everywhere with value
\(+\infty\) assigned to the singular factor.
Insert this inequality into \eqref{jk-after-scaling}.
Proposition~\ref{jk-weighted} bounds the resulting integral of
\(F(x)|b_0|^{-\gamma}\).
Together with \eqref{jk-after-rotation}, this proves
\eqref{jk-os}.
\end{proof}

\subsection{Sharp convergence restrictions and scope}

The inequalities in \eqref{jk-cone} are necessary, not merely
convenient sufficient assumptions.  The proof of
Theorem~\ref{jk-exact} gives a divergent complex coordinate integral
or the divergent final real or radial integral whenever any of
them fails.  For example, \(p=0,\nu_0=0\) gives logarithmic
divergence; for \(p=1\), the vectors
\((0,1)\) and \((2,-1)\) respectively violate the first
complex integration condition and the second one at equality,
although the remaining inequalities are positive.

The extra weight in Theorem~\ref{jk-oscillatory} is also essential.
For \(P\geq1\) and \(e_0>0\), the integral of
\(\prod_l d_l^{-\sigma_l+e_l}\) without that weight is infinite,
because its first exponent parameter is \(\nu_0=-e_0<0\).
If all the coefficients \(\eta,\eta',\xi\) vanish, the weight
is a positive constant, so the same divergence persists.
This supplies a counterexample to omitting the condition on the
linear form; it is not a counterexample under the stated
condition \(\eta'_0\ne0\).
In Lemma~\ref{jk-riesz}, \(\gamma\geq2\) makes the origin
nonintegrable, and \(\gamma\leq e\) makes infinity nonintegrable.
The restriction \(e>0\) is also required for a bound uniform
as \(D\downarrow0\) with \(c\ne0\): a disc around \(c\)
has a divergent limiting integral when \(e\leq0\).

The definitions \eqref{jk-coordinates} and the critical exponents
are exactly those in the source's coordinate integral
\(\mathrm{CHInt}\).  Theorem~\ref{jk-exact} proves its assertion
\(\mathrm{CH}\) and strengthens it to an exact evaluation and a
necessary convergence criterion.  Corollary~\ref{jk-tail}
proves its assertion \(\mathrm{TailFin}\), identity
\eqref{jk-square} is its \(\mathrm{C1}\), and the three elementary
lemmas prove \(\mathrm{SUB}\), \(\mathrm{KER}\), and
\(\mathrm{SLOW}\).  Theorem~\ref{jk-oscillatory}, with \(P=p+1\),
is exactly its assertion \(\mathrm{OS}\).
These identifications refer to the definitions in
\texttt{TailIntegrals/TailInt.lean} in the source archive
\cite{alpo-uniform-source}; the accompanying source trace records
the other files and the parts read.

This section supplies actual estimates, with proofs, rather than
assuming the existence of a Jacquet majorant.
To obtain that majorant one still has to prove the geometric
identities expressing the unipotent integral in these coordinates,
the corresponding induction on the rank, and the group integral
used for the matrix coefficient bound.
The convergence and differentiability of the final Jacquet
pairing also require the subsequent unfolding and domination
arguments.  None of those further assertions is inferred solely
from the coordinate estimates proved here.


